%% slide07_study_note.tex
%% Detailed Study Note --- Slide 7: Ch 4 SAMBP-87L Adaptive Line Differential
%% TSA Meeting Preparation | Anoop V Eluvathingal | IIT Madras | 2026
\documentclass[12pt,a4paper]{article}

\usepackage[margin=2.5cm]{geometry}
\usepackage{amsmath,amssymb,bm}
\usepackage{booktabs,tabularx,array,longtable}
\usepackage{graphicx}
\usepackage{xcolor}
\usepackage{tcolorbox}
\usepackage{enumitem}
\usepackage{fancyhdr}
\usepackage{titlesec}
\usepackage{hyperref}
\usepackage{parskip}
\usepackage{algorithm}
\usepackage{algpseudocode}
\algrenewcommand\algorithmicrequire{\textbf{Input:}}
\algrenewcommand\algorithmicensure{\textbf{Output:}}
\newcommand{\kibr}{k_\mathrm{ibr}}
\newcommand{\Idiff}{I_\mathrm{diff}}
\newcommand{\Ibias}{I_\mathrm{bias}}
\newcommand{\km}{k_m}
\newcommand{\kbase}{k_\mathrm{base}}
\newcommand{\kfloor}{k_\mathrm{floor}}

\definecolor{iitmnavy}{RGB}{15,40,80}
\definecolor{iitmgold}{RGB}{180,140,0}
\definecolor{lightblue}{RGB}{220,235,250}
\definecolor{lightyellow}{RGB}{255,252,220}
\definecolor{lightred}{RGB}{255,235,235}
\definecolor{lightgreen}{RGB}{230,255,230}
\definecolor{lightpurple}{RGB}{240,230,255}
\definecolor{lightorange}{RGB}{255,243,225}

\tcbuselibrary{skins,breakable}
\newtcolorbox{keybox}[1]{colback=lightblue,colframe=iitmnavy,
  fonttitle=\bfseries\color{white},title=#1,arc=3pt,boxrule=1pt,breakable}
\newtcolorbox{formulabox}[1]{colback=lightyellow,colframe=iitmgold,
  fonttitle=\bfseries,title=#1,arc=3pt,boxrule=1pt,breakable}
\newtcolorbox{resultbox}[1]{colback=lightgreen,colframe=green!50!black,
  fonttitle=\bfseries,title=#1,arc=3pt,boxrule=1pt,breakable}
\newtcolorbox{warnbox}[1]{colback=lightred,colframe=red!60!black,
  fonttitle=\bfseries,title=#1,arc=3pt,boxrule=1pt,breakable}
\newtcolorbox{archbox}[1]{colback=lightpurple,colframe=iitmnavy!60,
  fonttitle=\bfseries,title=#1,arc=3pt,boxrule=1pt,breakable}
\newtcolorbox{speakbox}{colback=orange!8,colframe=orange!70!black,
  fonttitle=\bfseries,title={What to Say at the TSA Meeting},arc=3pt,
  boxrule=1pt,breakable}

\pagestyle{fancy}
\fancyhf{}
\fancyhead[L]{\color{iitmnavy}\small\textbf{TSA Meeting Study Note}}
\fancyhead[R]{\color{iitmnavy}\small Slide 7 --- Ch~3 SAMBP-87L Adaptive Line Differential}
\fancyfoot[C]{\thepage}
\fancyfoot[R]{\small\color{gray}Anoop V Eluvathingal $\cdot$ IIT Madras $\cdot$ 2026}
\renewcommand{\headrulewidth}{1.5pt}
\renewcommand{\headrule}{\color{iitmnavy}\hrule width\headwidth height\headrulewidth}

\titleformat{\section}{\large\bfseries\color{iitmnavy}}{}{0em}{}[\color{iitmnavy}\titlerule]
\titleformat{\subsection}{\normalsize\bfseries\color{iitmnavy!80}}{}{0em}{}

\hypersetup{colorlinks=true,linkcolor=iitmnavy,urlcolor=iitmnavy}

\begin{document}

%% ── Title Block ──────────────────────────────────────────────────────────────
\begin{tcolorbox}[colback=iitmnavy,coltext=white,arc=4pt,boxrule=0pt]
\begin{center}
  {\Large\bfseries TSA Meeting --- Slide 7 Study Note}\\[4pt]
  {\large Ch~3: SAMBP-87L --- Adaptive Line Current Differential Protection}\\[6pt]
  {\normalsize Contributions C1, C2, C3, C4 $\cdot$ Chapter 3 $\cdot$ SAMBP Framework}\\[2pt]
  {\small Anoop V Eluvathingal $\cdot$ IIT Madras $\cdot$ PhD Thesis 2026}
\end{center}
\end{tcolorbox}

\vspace{6pt}

%% ── 1. Slide Overview ────────────────────────────────────────────────────────
\section{Slide Overview and Narrative}

\begin{keybox}{What This Slide Shows}
Slide~7 presents \textbf{Ch~3: SAMBP-87L}, the adaptive line current differential relay. The slide shows two complementary aspects:

\textbf{Left panel --- Adaptive slope law:} The central equation is the IBR-aware bias slope:
\[
  \km(\kibr) = \kbase \cdot \frac{(1-\kibr)^2}{(1+\kibr)^2} + \kfloor
\]
with two limiting behaviours: as $\kibr \to 1$ (IBR dominates), $\km \to \kfloor$ (minimum slope, preserving sensitivity); as $\kibr \to 0$ (SG network), $\km \to \kbase$ (classical conservative setting). The Stage-2 veto gate requires $\gamma \geq 0.85$ before the trip command is issued. The Lyapunov stability proof --- $V = \|\kibr - \hat{\kibr}\|^2$, $\dot{V} < 0$ under bounded CT noise --- guarantees estimation convergence.

\textbf{Right panel --- Relay logic:} Six-step execution loop: (1) sample phasors $I_S$, $I_R$ at 1~kHz; (2) compute $\kibr$ via LM inverse estimator; (3) update $\km(\hat{\kibr})$ every 20~ms; (4) evaluate the operate/restrain criterion; (5) Stage-2 veto (check $\gamma \geq 0.85$); (6) issue trip if both stages pass.

\textbf{Key observation:} The 20~ms update latency is within the IEC~60255 protection-response window, meaning the adaptive boundary tracks IBR operating changes fast enough to remain valid throughout the fault transient.
\end{keybox}

\subsection{Four Contributions in Chapter~3}

The slide focuses on C1 (adaptive slope). The full chapter delivers four complementary contributions:

\begin{center}
\begin{tabular}{@{}clll@{}}
\toprule
\textbf{ID} & \textbf{Element} & \textbf{Target fault type} & \textbf{Key mechanism} \\
\midrule
C1 & Adaptive 87L  & All fault types, main element & Adaptive slope $\km(\kibr)$ \\
C2 & 87LH harmonic & IBR harmonic, blocks inrush  & 2nd/5th harmonic ratio $h_2 < 0.15$ \\
C3 & 87LN neg-seq  & Asymmetric (SLG, LL, LLG)   & $\rho_2 = 2$ guarantee \\
C4 & TDCS-87L      & 3-phase faults on IBR-only lines & Transient $dI/dt$ spike \\
\bottomrule
\end{tabular}
\end{center}

\textbf{Combined result (TR-23/TR-27):} Baseline 87L alone covers 12/21 scenarios (57\%). Adding 87LN (C3) recovers 5 SLG cases: 17/21 (81\%). Adding TDCS-87L (C4) recovers the remaining 4 symmetric 3-phase IBR scenarios: 21/21 (100\%).

%% ── 2. Engineering Challenge ─────────────────────────────────────────────────
\section{Engineering Challenge: Why Fixed Slope Fails at High IBR Penetration}

\begin{warnbox}{The Core Problem: Operate-to-Bias Ratio Collapses with IBR}
The classical 87L relay operates when $\Idiff > \km \cdot \Ibias + I_0$. During an internal fault:
\[
  \Idiff = |\hat{I}_A + \hat{I}_B|, \qquad \Ibias = \frac{|\hat{I}_A| + |\hat{I}_B|}{2}
\]
where $\hat{I}_A$ is the grid-fed current and $\hat{I}_B = \kibr \cdot I_\mathrm{rated}$ is the IBR current. The operate-to-bias ratio (OBR) is:
\begin{equation}
  \rho = \frac{\Idiff}{\Ibias} = \frac{2\sqrt{I_A^2 + I_B^2 + 2I_A I_B \cos\Delta\phi}}{I_A + I_B}
  \label{eq:rho_general}
\end{equation}
The \emph{minimum} OBR (worst case: sources in anti-phase, $\cos\Delta\phi = -1$) is:
\begin{equation}
  \rho_\mathrm{min} = \frac{2|I_A - I_B|}{I_A + I_B}
    = \frac{2(1-\kibr)}{1+\kibr}
  \label{eq:rho_min}
\end{equation}

\medskip
\begin{tabular}{@{}lrll@{}}
\toprule
$\kibr$ & $\rho_\mathrm{min}$ & Fixed slope $\km = 0.30$ & Verdict \\
\midrule
0.10 & 1.636 & $1.636 > 0.30$ & \textbf{Operate} (easy) \\
0.40 & 0.857 & $0.857 > 0.30$ & \textbf{Operate} \\
0.50 & 0.667 & $0.667 > 0.30$ & \textbf{Operate} (marginal) \\
0.70 & 0.176 & $0.176 < 0.30$ & \textbf{FAIL --- no trip} \\
0.85 & 0.081 & $0.081 < 0.30$ & \textbf{FAIL --- no trip} \\
\bottomrule
\end{tabular}

\medskip
At $\kibr = 0.70$, the minimum OBR ($\rho_\mathrm{min} = 0.176$) falls below the fixed slope ($0.30$). The relay does not trip even for a solid internal fault. The solution: make $\km$ decrease as $\kibr$ increases. But decrease it \emph{safely} --- not so far that CT errors cause a false trip.
\end{warnbox}

\subsection{The Fundamental Tension}

There is a design tension between two requirements:
\begin{itemize}[nosep]
  \item \textbf{Sensitivity} (trip for internal faults): requires $\km < \rho_\mathrm{min} = 2(1-\kibr)/(1+\kibr)$.
  \item \textbf{Stability} (no false trip during CT errors): requires $\km > 2\epsilon_\mathrm{CT}$ (Class~1P: $\km > 0.02$).
\end{itemize}
A fixed slope cannot satisfy both simultaneously across all $\kibr$: for low $\kibr$, $\km = 0.30$ provides a 15:1 stability margin but is unnecessarily conservative for sensitivity; for high $\kibr$, $\km = 0.30$ provides adequate stability but exceeds $\rho_\mathrm{min}$, blocking the trip. The adaptive slope resolves this tension by tracking $\kibr$ in real time.

%% ── 3. Adaptive Slope Derivation ─────────────────────────────────────────────
\section{Adaptive Slope Derivation (Contribution C1)}

\begin{formulabox}{Adaptive Bias Slope: Physics Derivation}
The optimal slope simultaneously satisfies stability and sensitivity at every $\kibr$:
\begin{equation}
  \km^\ast(\kibr) = \max\!\left[
    \underbrace{2\epsilon_\mathrm{CT} + \frac{I_c}{I_\mathrm{bias}}}_{\text{stability floor}},\;
    \frac{1}{1-\delta_m}\underbrace{\frac{2\sqrt{1+\kibr^2-2\kibr\cos\psi}}{1+\kibr}}_{\text{sensitivity ceiling}}
  \right]
  \label{eq:km_optimal}
\end{equation}
where $\psi = \arg(Z_F/(Z_S + Z_L)) + \pi$ is the fault angle and $\delta_m = 0.05$ (5\% design margin).

\textbf{Simplified closed form (the slide equation):}
\begin{equation}
  \km(\kibr) = \kbase \cdot \frac{(1-\kibr)^2}{(1+\kibr)^2} + \kfloor
  \label{eq:km_slide}
\end{equation}
where $\kbase \approx 0.5$ and $\kfloor \approx 0.05$.

\textbf{Physical interpretation of the two limits:}
\begin{itemize}[nosep]
  \item $\kibr \to 0$ (all-SG network): $(1-0)^2/(1+0)^2 = 1$, so $\km \to \kbase + \kfloor \approx 0.55$. This is the conservative high-slope setting used in classical SG-only networks.
  \item $\kibr \to 1$ (IBR-dominated): $(1-1)^2/(1+1)^2 = 0$, so $\km \to \kfloor \approx 0.05$. The slope reaches its minimum --- just above the CT error floor --- to preserve sensitivity for the weak IBR fault current.
\end{itemize}

\textbf{Why $(1-\kibr)^2/(1+\kibr)^2$?} This ratio tracks $\rho_\mathrm{min}^2/4$, the squared operate-to-bias ratio at worst-case phase angle. The adaptive slope therefore stays always just below the sensitivity ceiling while remaining above the stability floor.
\end{formulabox}

\subsection{Real-Time k-ibr Estimation}

The relay estimates $\kibr$ from the received remote current phasor:
\begin{equation}
  \hat{\kibr}(t) = \frac{|\hat{I}_R(t)|}{I_\mathrm{rated}}
  \label{eq:kibr_est}
\end{equation}
where $\hat{I}_R$ is received over IEC~61850 GOOSE within 4~ms from the remote relay ($I_R \equiv I_B$ in the slide notation, Remote end current). A first-order smoothing filter prevents noise spikes from causing erratic slope updates:
\begin{equation}
  \hat{\kibr}(t) \leftarrow (1-\lambda)\hat{\kibr}(t-1) + \lambda\hat{\kibr}_\mathrm{raw}(t), \quad \lambda = 0.10
  \label{eq:kibr_filter}
\end{equation}
Time constant: $\approx 10$ samples at 1~kHz sampling = 10~ms. The slope is then updated every 20~ms (two filter time constants), ensuring the estimate has converged before it is applied.

%% ── 4. Lyapunov Stability Proof ──────────────────────────────────────────────
\section{Lyapunov Stability Proof}

\begin{archbox}{The Slide Claim: V-dot less than zero under bounded CT noise}
The slide states a Lyapunov stability result for the $\kibr$ estimator:
\[
  V = \|\kibr - \hat{\kibr}\|^2, \qquad \dot{V} < 0 \;\text{under bounded CT noise}
\]

\textbf{What this means:} $V$ is the squared estimation error --- the gap between the true IBR penetration $\kibr$ and the relay's real-time estimate $\hat{\kibr}$. The Lyapunov condition $\dot{V} < 0$ (negative time derivative) means the estimation error is strictly decreasing --- the estimator converges. This holds as long as the CT measurement noise is bounded (not arbitrarily large).

\textbf{Proof sketch:}
\begin{enumerate}[nosep]
  \item Define the estimation update: $\hat{\kibr}_{k+1} = (1-\lambda)\hat{\kibr}_k + \lambda \cdot (|\hat{I}_R|/I_\mathrm{rated})$.
  \item The CT-corrupted measurement is $|\hat{I}_R| = \kibr I_\mathrm{rated} + \nu$ where $|\nu| \leq \bar{\nu}$ (bounded noise).
  \item The update error: $\kibr - \hat{\kibr}_{k+1} = (1-\lambda)(\kibr - \hat{\kibr}_k) - \lambda\nu$.
  \item Then $V_{k+1} = (1-\lambda)^2 V_k + \lambda^2\nu^2 - 2\lambda(1-\lambda)(\kibr-\hat{\kibr}_k)\nu$.
  \item In expectation (zero-mean noise): $\mathbb{E}[V_{k+1}] = (1-\lambda)^2 \mathbb{E}[V_k] + \lambda^2\bar{\nu}^2$.
  \item For $(1-\lambda)^2 < 1$ (satisfied for $\lambda \in (0,1)$): $\mathbb{E}[V_{k+1}] < \mathbb{E}[V_k]$ whenever $V_k > \lambda^2\bar{\nu}^2/(1-(1-\lambda)^2)$ --- i.e., the estimator converges to an error ball of radius proportional to $\bar{\nu}$ (bounded by CT noise).
\end{enumerate}

\textbf{Why it matters to the examiner:} This proves the adaptive relay is \emph{not} ``guessing'' the IBR level --- it provably converges to the correct $\kibr$ and therefore applies the correct slope. Without this proof, the examiners could argue that $\hat{\kibr}$ might diverge or oscillate.
\end{archbox}

%% ── 5. Four Elements: C2, C3, C4 ────────────────────────────────────────────
\section{The Three Additional Elements: C2, C3, C4}

\subsection{C2: Harmonic Discrimination 87LH}

IBR fault current contains 2nd and 5th harmonic components from filter resonance during current controller transients. The harmonic ratio is:
\begin{equation}
  h_2 = \frac{I_2}{I_1} \approx 0.15 \quad \text{(GFL inverter, 500~Hz bandwidth)}
  \label{eq:h2}
\end{equation}
The 87LH element trips when the fundamental differential exceeds threshold ($I_{\mathrm{diff},1} > \km \Ibias + I_0$) \textbf{AND} the harmonic ratio is below the inrush blocking threshold:
\begin{equation}
  \frac{I_{\mathrm{diff},2}}{I_{\mathrm{diff},1}} < h_\mathrm{block} = 0.15
  \label{eq:87lh}
\end{equation}
This specifically distinguishes IBR faults ($h_2 \approx 0.15$, large fundamental differential) from transformer inrush ($h_2 > 0.20$, near-zero fundamental differential). Note: \textbf{the harmonic block threshold is lower than the standard inrush criterion of 0.20} --- providing a 5\% margin that prevents the 87LH from blocking on IBR fault currents.

\subsection{C3: Negative-Sequence 87LN Element}

\begin{resultbox}{Key Result: IBR Suppression Helps 87LN (not hurts)}
For asymmetric faults (SLG, LL, LLG), IBR negative-sequence suppression produces a counterintuitive benefit:

When IBR suppresses its negative-sequence output ($\hat{I}_{B,2}^\mathrm{IBR} = 0$), the negative-sequence differential current equals only the grid-side contribution:
\[
  \hat{I}_{\mathrm{diff},2} = \hat{I}_{A,2}
\]
The negative-sequence bias is $I_{\mathrm{bias},2} = (I_{A,2} + 0)/2 = I_{A,2}/2$.

Therefore the negative-sequence OBR is:
\begin{equation}
  \rho_2 = \frac{\hat{I}_{\mathrm{diff},2}}{I_{\mathrm{bias},2}} = \frac{I_{A,2}}{I_{A,2}/2} = 2
  \label{eq:rho2}
\end{equation}

$\rho_2 = 2$ \textbf{regardless of $\kibr$}. Since the slope threshold $k_{m,2} = 1.5 < 2$, the 87LN element \textit{always} operates for any asymmetric internal fault when the IBR suppresses negative-sequence. IBR sequence suppression --- usually viewed as a weakness for protection --- here becomes an \textbf{asset}.
\end{resultbox}

\subsection{C4: Transient Differential Current Signature (TDCS-87L)}

Three-phase IBR-only faults produce $I_2 = I_0 = 0$ at the relay terminal --- making 87LN and all sequence elements inactive. TDCS-87L exploits the travelling-wave transient in the first 5~ms post-fault:

\begin{equation}
  S_\mathrm{TDCS}(t) = \sum_{k=t-N_w}^{t}
  \left[\frac{d}{dt}\bigl(\Idiff(t)\bigr)\right]^2 \Delta t
  \label{eq:tdcs}
\end{equation}
where $N_w = 5$~ms sliding window. An internal fault produces a sharp $dI_\mathrm{diff}/dt$ spike when the travelling wave arrives. External faults and load steps produce near-zero differential transients (Kirchhoff: $\Delta I = 0$ for balanced load steps). Threshold: $S_\mathrm{TDCS} > 3\sigma$ above pre-fault RMS.

\textbf{TR-23 results:} TDCS detection limit $k_\mathrm{det,min} \geq 0.040$~pu. Security: 7/7 balanced load steps correctly rejected (no false trips). External fault security: 7/7 correct (CT mismatch $\leq 0.5\%$, margin $1.77\times$).

%% ── 6. Stage-2 Veto Gate ─────────────────────────────────────────────────────
\section{Stage-2 Veto Gate}

\begin{archbox}{Stage-2 Gate: gamma greater than or equal to 0.85}
The Stage-2 veto gate is the SAMBP framework's quality supervisor. For the 87L relay, the gate requires $\gamma \geq 0.85$ before any trip command is issued (higher threshold than the 0.70 used in SAMBP-SyncOC, reflecting the greater consequence of a line differential mis-trip).

$\gamma$ is computed from the LM estimator quality metrics:
\begin{itemize}[nosep]
  \item \textbf{$\gamma_R$} (residual fit): $R^2$ of the current waveform fit to the differential model.
  \item \textbf{$\gamma_J$} (Jacobian condition): $\exp(-\alpha_J \kappa_n)$; poor conditioning (CT saturation) lowers this.
  \item \textbf{$\gamma_\Delta$} (stability): change in $\hat{\kibr}$ from the previous estimate.
  \item \textbf{$\gamma_P$} (physical bounds): all parameters within the feasible set $\Theta$.
\end{itemize}

\textbf{If $\gamma < 0.85$:} Trip is vetoed. The relay holds its current decision. This prevents spurious trips during:
\begin{itemize}[nosep]
  \item CT saturation (high-load through-fault): $\kappa_n$ spikes, $\gamma_J$ drops.
  \item Line energisation transients: large $\hat{\kibr}$ change, $\gamma_\Delta$ drops.
  \item Communication interruption (GOOSE loss): $\hat{I}_R$ unavailable, $\gamma_P$ fails.
\end{itemize}

\textbf{Timing:} Stage-2 check adds $< 1$~ms to the relay cycle (confidence score is evaluated in the same DSP loop as the LM estimator). Total relay cycle: phasor sampling (1~ms) + LM estimation (5~ms) + Stage-2 check ($< 1$~ms) + slope update (1~ms) = $\approx 8$~ms per trip decision.
\end{archbox}

%% ── 7. Results ───────────────────────────────────────────────────────────────
\section{Results: Cumulative Coverage from 57 to 100 Percent}

\begin{resultbox}{TR-23 Figure 3: Cascading Improvement Across Three SAMBP Generations}
The 21-scenario test suite covers the full $(\kibr, R_F, \text{fault type})$ operating space. The three bars show cumulative coverage as each element is added:

\medskip
\begin{center}
\begin{tabular}{@{}clrll@{}}
\toprule
\textbf{Generation} & \textbf{Configuration} & \textbf{Scenarios correct} & \textbf{Coverage} & \textbf{Gap closed} \\
\midrule
Baseline & 87L only (fixed slope)  & 12/21 & 57\% & --- \\
TR-22    & 87L + 87LN              & 17/21 & 81\% & SLG with IBR neg-seq suppr. \\
TR-23    & 87L + 87LN + TDCS       & 21/21 & 100\% & Symmetric 3PH, IBR-only lines \\
\bottomrule
\end{tabular}
\end{center}

\medskip
\textbf{The 9 baseline failures:}
\begin{itemize}[nosep]
  \item 5 SLG failures: $\kibr < 0.12$~pu --- baseline 87L fundamental current below threshold; recovered by 87LN.
  \item 4 three-phase failures: symmetric fault on IBR-only line --- $I_2 = I_0 = 0$, no sequence element active; recovered by TDCS.
\end{itemize}

\textbf{TR-27 integration test (21 scenarios including IBR-specific):} All 21/21 correctly resolved by the combined protection chain (87L + 87LN + TDCS + adaptive slope). No false trips in 7 load-step security tests.
\end{resultbox}

%% ── 8. Best Figure Recommendation ───────────────────────────────────────────
\section{Best Figure to Show at the TSA Meeting}

\begin{keybox}{Recommended Figure: TR-23 Figure 3 --- Cumulative Scenario Coverage Bar Chart}
\textbf{What it shows:} Three grouped bars (Baseline 87L only, TR-22 with 87LN, TR-23 with TDCS) with colour coding: dark green = baseline correct, medium green = recovered by 87LN (SLG), light green = recovered by TDCS (3PH IBR), red = undetected failure. The chart shows the progression 12/21 (57\%) $\to$ 17/21 (81\%) $\to$ 21/21 (100\%).

\textbf{Why it is the best figure:}
\begin{enumerate}[nosep]
  \item It is the thesis's ``zero-gap'' result in visual form --- a single figure proves all 21 scenarios are covered.
  \item Each colour segment identifies which element closed which gap, making the design logic immediately visible.
  \item Examiners can immediately see that the 4 TDCS-recovered cases (light green) would have been completely missed without C4, and can ask an intelligent question about why.
  \item It does not require any mathematical knowledge to understand --- the story is told entirely visually.
\end{enumerate}

\textbf{Where to find it:} TR-23/2026, page~5, Figure~3. Source: SAMBP test suite, integration module.

\textbf{Secondary figure:} The adaptive slope curve $\km^\ast(\kibr)$ vs $\kibr$ (ch04 Figure~1). This is the best figure for explaining the slide's main equation: it shows the classical fixed slope (dashed line at $\km = 0.30$) crossing the sensitivity ceiling for $\kibr > 0.5$, and the adaptive slope (solid curve) tracking just below the ceiling while staying above the stability floor. Draw this on the whiteboard if the figure is not printed.

\textbf{Whiteboard supplement (60 seconds):} Draw two curves vs $\kibr$: (1) the sensitivity ceiling $\rho_\mathrm{min} = 2(1-\kibr)/(1+\kibr)$ --- a decreasing curve from 2 at $\kibr=0$ to 0 at $\kibr=1$; (2) a horizontal fixed slope at 0.30. Mark where they cross ($\kibr \approx 0.54$) as the ``over-restraint boundary''. Then draw the adaptive slope tracking below the ceiling. This makes the design motivation unmistakable.
\end{keybox}

%% ── 9. Cross-Thesis Connections ──────────────────────────────────────────────
\section{Cross-Thesis Connections}

\begin{archbox}{How This Chapter Connects to Other Work}
\textbf{Depends on (upstream):}
\begin{itemize}[nosep]
  \item \textbf{Ch~2 (SAMBP Foundations)}: The LM inverse estimator, Stage-2 gate, Lyapunov stability framework, and bounded update law are all developed in Ch~2 and applied here. This is the \emph{first thesis chapter to deploy the full SAMBP framework}.
  \item \textbf{Ch~3 (IBR Characterisation)}: The failure modes addressed --- fixed slope insensitive for $\kibr > 0.5$, sequence element blind spots on IBR lines --- are identified as F2 and F3 in Ch~3.
\end{itemize}

\textbf{Enables (downstream):}
\begin{itemize}[nosep]
  \item \textbf{Ch~6 (SAMBP-87T)}: The adaptive slope concept is extended from line to transformer differential in Ch~6. The SPRT replaces the threshold comparison, but the same $\hat{\kibr}$-based adaptation drives the likelihood model.
  \item \textbf{Ch~5 (Distance Protection)}: The TDCS concept (transient differential current) is adapted in Ch~7 to the Zone~1 adaptive reach update. Both exploit the travelling-wave transient as a discriminating feature.
  \item \textbf{Ch~7 (WAPC)}: The 87L trip signal feeds into the Wide-Area Protection Coordinator's N-1 logic. The WAPC uses 87L as the ``indispensable back-stop'' for all $\kibr \geq 0.06$ scenarios (TR-35 terminology).
  \item \textbf{Ch~8 (Cybersecurity)}: The GOOSE messages carrying $\hat{I}_R$ from the remote relay are the attack surface for GOOSE spoofing and replay attacks (Ch~12/TR-52). HMAC-SHA256 authentication adds 0.36~ms overhead --- negligible against the 20~ms update period.
\end{itemize}

\textbf{Design pattern across SAMBP:} Ch~4 establishes the 3-layer pattern used in all subsequent protection chapters: (1) Physics-derived adaptation law; (2) Stage-2 quality gate; (3) Hierarchical vote of specialised elements. Chapters 4, 5, 6, 7, 8 all use this same architecture.
\end{archbox}

%% ── 10. Key Equations Reference Card ────────────────────────────────────────
\section{Key Equations Reference Card}

\begin{formulabox}{Equations You Must Be Able to Write and Explain}

\textbf{1. Classical 87L criterion:}
\[
  \Idiff > \km \cdot \Ibias + I_0, \quad
  \Idiff = |\hat{I}_S + \hat{I}_R|, \quad
  \Ibias = \tfrac{|\hat{I}_S| + |\hat{I}_R|}{2}
\]
Say: ``Differential = phasor sum of send and receive currents. Bias = average magnitude. The criterion is: differential exceeds slope times bias plus a minimum pickup.''

\textbf{2. Minimum operate-to-bias ratio (why the problem exists):}
\[
  \rho_\mathrm{min} = \frac{2(1-\kibr)}{1+\kibr}
\]
Say: ``At $\kibr = 0.7$, $\rho_\mathrm{min} = 0.18$. A fixed slope of 0.30 blocks the trip. The relay is over-restrained. This is the core problem.''

\textbf{3. Adaptive slope (slide equation):}
\[
  \km(\kibr) = \kbase \cdot \frac{(1-\kibr)^2}{(1+\kibr)^2} + \kfloor
\]
Say: ``$k_\mathrm{base} \approx 0.5$ tracks the sensitivity ceiling. $k_\mathrm{floor} \approx 0.05$ is the minimum slope, set above the CT error floor. As $k_\mathrm{ibr} \to 1$, only the floor remains.''

\textbf{4. 87LN negative-sequence OBR (the key insight):}
\[
  \rho_2 = \frac{\hat{I}_{\mathrm{diff},2}}{I_{\mathrm{bias},2}}
    = \frac{I_{A,2}}{I_{A,2}/2} = 2, \quad \text{when } I_{B,2}^\mathrm{IBR} = 0
\]
Say: ``When the IBR suppresses its negative-sequence output --- which it does for grid compliance --- the OBR is exactly 2 regardless of $k_\mathrm{ibr}$. IBR suppression makes 87LN more sensitive, not less.''

\textbf{5. TDCS score:}
\[
  S_\mathrm{TDCS}(t) = \sum_{k=t-N_w}^{t} \!\left[\frac{d\Idiff}{dt}\right]^2 \Delta t, \quad N_w = 5\;\text{ms}
\]
Say: ``This is the energy of the differential current derivative in a 5~ms window. An internal fault produces a sharp travelling-wave spike. A balanced load step produces zero differential current by Kirchhoff, so $S_\mathrm{TDCS} = 0$.''

\textbf{6. Lyapunov stability:}
\[
  V = \|\kibr - \hat{\kibr}\|^2, \qquad \dot{V} < 0 \;\text{under bounded CT noise}
\]
Say: ``$V$ is the estimation error squared. $\dot{V} < 0$ means the error decreases --- the estimator converges. The proof uses the exponential filter's contraction property. The adaptive slope is therefore tracking a convergent estimate.''
\end{formulabox}

%% ── 11. Speaker Script ───────────────────────────────────────────────────────
\section{Speaker Script (Timed to 3 Minutes)}

\begin{speakbox}
\textbf{[0:00--0:35] Open with the problem:}\\
``Line current differential protection is the gold standard for line protection: it needs no voltage measurement, it is immune to power swings, and it detects high-impedance faults that defeat every other protection scheme. But it has one critical vulnerability with IBR: the bias slope. The classical relay operates when the differential current exceeds a fixed percentage of the bias current. As IBR penetration increases, the IBR current-limits to around 1 pu of rated during a fault. The grid side still sees the fault current, but the IBR side contribution is small. The operate-to-bias ratio collapses as $2(1-k_\mathrm{ibr})/(1+k_\mathrm{ibr})$. At $k_\mathrm{ibr} = 0.7$, that ratio is 0.18 --- less than the fixed slope of 0.30. The relay refuses to trip. That is over-restraint.''

\textbf{[0:35--1:20] Explain the adaptive slope:}\\
``SAMBP-87L fixes this with an adaptive slope $k_m(k_\mathrm{ibr}) = k_\mathrm{base} \cdot (1-k_\mathrm{ibr})^2/(1+k_\mathrm{ibr})^2 + k_\mathrm{floor}$. As $k_\mathrm{ibr}$ increases toward 1, the slope automatically decreases toward $k_\mathrm{floor} = 0.05$ --- the minimum value that keeps the relay stable against CT measurement error. As $k_\mathrm{ibr}$ decreases toward 0, the slope increases to the classical conservative setting. The relay estimates $k_\mathrm{ibr}$ from the received remote current phasor every 20~ms. A Lyapunov stability proof --- $V = ||k_\mathrm{ibr} - \hat{k}_\mathrm{ibr}||^2$, $\dot{V} < 0$ --- guarantees the estimate converges under bounded CT noise. Before issuing any trip, the Stage-2 veto gate checks $\gamma \geq 0.85$ to confirm the estimate is trustworthy.''

\textbf{[1:20--2:00] Explain the three additional elements:}\\
``The adaptive slope handles most scenarios, but three gaps remain. For SLG faults when $k_\mathrm{ibr}$ is very low --- below 0.12 --- the fundamental differential current is too small to trigger the main element. The 87LN negative-sequence element fills this gap. Here is a key insight: when the IBR suppresses its negative-sequence output for grid compliance, the negative-sequence OBR is always exactly 2, regardless of $k_\mathrm{ibr}$. IBR suppression actually helps 87LN. The second gap is symmetric three-phase faults on IBR-only lines, where all sequence elements are blind. The TDCS element uses the differential current derivative in a 5~ms window to detect the travelling-wave spike from the fault inception. Balanced load steps produce zero differential current by Kirchhoff's law, so TDCS is immune to them.''

\textbf{[2:00--2:40] Present the results:}\\
``The TR-23 cumulative coverage chart shows the progression clearly. The baseline 87L alone correctly handles 12 of 21 test scenarios --- 57\%. Adding the 87LN element recovers 5 SLG cases, taking coverage to 17/21 or 81\%. Adding TDCS recovers the remaining 4 symmetric three-phase IBR scenarios: 21/21, 100\% coverage. The system integration test in TR-27 confirms all 21 scenarios resolved with no false trips on 7 balanced load-step security tests.''

\textbf{[2:40--3:00] Close:}\\
``This is Contributions C1 through C4. Each element is physics-derived: the adaptive slope from the sensitivity-stability constraint, 87LN from sequence suppression analysis, TDCS from travelling-wave physics. Together they cover the full IBR operating envelope that the classical 87L cannot reach. The 20~ms update latency is within the IEC~60255 protection-response window, confirmed by measurement.''
\end{speakbox}

%% ── 12. Anticipated Q and A ──────────────────────────────────────────────────
\section{Anticipated Questions and Answers}

\begin{keybox}{Q1: Could the adaptive slope go too low and cause a false trip during through-fault CT saturation?}
\textbf{Answer:} This is the core security concern, and it is directly addressed by the stability condition. The adaptive slope has a hard floor $\kfloor = 0.05$, which is above $2\epsilon_\mathrm{CT} = 0.02$ for Class~1P CTs. So even at $\kibr = 1$ (minimum slope), the relay is stable for CT errors up to 2.5\%. For Class~1P CTs ($\epsilon_\mathrm{CT} \leq 1\%$), the stability margin is $0.05 - 0.02 = 0.03$ --- a 50\% buffer above the CT error. The Proposition (stability under CT saturation) in the chapter proves this formally: provided the CT composite error satisfies $\epsilon_\mathrm{CT} \leq (k_m^*(\kibr) - I_c/I_\mathrm{bias})/2$, the relay never false-trips on a through-fault. Additionally, the Stage-2 gate will veto any trip where the Jacobian condition number (indicating CT saturation) is elevated.
\end{keybox}

\begin{keybox}{Q2: Why does the 87LN use k-m2 = 1.5 when the guaranteed OBR is 2? Why not use 1.9?}
\textbf{Answer:} Setting $k_{m,2} = 1.5$ provides a 25\% margin below the guaranteed $\rho_2 = 2$. This margin accounts for: (1) partial IBR negative-sequence suppression (some IBRs inject limited $I_2$ for unbalance compensation under grid codes); (2) CT measurement error in the sequence extraction (unbalanced CT errors can corrupt $I_2$ estimation); (3) communication-induced phase error in $\hat{I}_{R,2}$. A setting of 1.9 would leave only a 5\% margin, insufficient for these practical error sources. The 1.5 setting is validated by TR-22: all 17 recovered SLG scenarios operate correctly with $k_{m,2} = 1.5$.
\end{keybox}

\begin{keybox}{Q3: The TDCS relies on a 5 ms window. What if the fault inception is missed in that window?}
\textbf{Answer:} The TDCS window is not a one-shot trigger. The score $S_\mathrm{TDCS}$ is computed on a \emph{sliding} 5~ms window, updated every sample (1~ms). If the travelling wave arrives at any point during the fault transient, the next 5~ms window will capture the $dI_\mathrm{diff}/dt$ spike. The relay sampling rate is 1~kHz, so the window advances every 1~ms and the TDCS evaluates every 1~ms. The first spike above $3\sigma$ threshold triggers the element. For a fault at $\alpha = 0.5$ on a 10~km line with wave speed $v_p = 0.8c$: $\tau_A = 5\,\mathrm{km}/(2.4\times10^8\,\mathrm{m/s}) \approx 20\,\mu\mathrm{s}$ --- the wave arrives within one sample. Detection latency is therefore 1--2~ms from fault inception.
\end{keybox}

\begin{keybox}{Q4: How is the relay logic different from conventional percentage-bias relays that use dual-slope characteristics?}
\textbf{Answer:} Conventional dual-slope relays use a two-slope bilinear characteristic: low slope (e.g., 0.25) for low bias current (through-load), high slope (e.g., 0.50) for high bias current (through-fault with CT saturation). The slope is a function of $I_\mathrm{bias}$, not $\kibr$. SAMBP-87L makes the slope a function of $\kibr$ --- the IBR penetration level --- which is a network characteristic, not a current magnitude. This is a fundamentally different parameterisation: a dual-slope relay would still use slope 0.30 for all $\kibr$ at a given $I_\mathrm{bias}$, causing over-restraint at $\kibr > 0.5$. The adaptive slope uses the same $k_m = 0.05$ at $\kibr = 0.8$ regardless of $I_\mathrm{bias}$, maintaining sensitivity where the OBR is inherently low.
\end{keybox}

\begin{keybox}{Q5: Does the 87LH threshold of 0.15 risk failing to block a real transformer inrush}
\textbf{Answer:} The 87LH element does not replace the transformer inrush blocking function --- it supplements the main 87L element for IBR-specific discrimination. The conventional 2nd-harmonic inrush block ($h_2 > 0.20$ $\Rightarrow$ block) is maintained separately in the SAMBP-87T (Ch~6) relay, which specifically handles transformer differential protection. The 87LH on the line differential side uses $h_\mathrm{block} = 0.15$ to prevent the harmonic ratio from being falsely elevated by IBR harmonic injection ($h_2 \approx 0.15$) from triggering an inrush-like block on a genuine line fault. The two elements cover different protected zones: 87LH covers the line, 87T (with SPRT) covers the transformer. A transformer inrush would not produce a differential current in the 87L measurement zone unless the transformer is inside the 87L zone, in which case the 87T operates first and the 87L is coordinated to reset.
\end{keybox}

\vspace{1em}
\begin{center}
\begin{tcolorbox}[colback=iitmnavy!10,colframe=iitmnavy,arc=4pt,boxrule=1pt,width=0.88\textwidth]
\centering
\textbf{Chapter~4 One-Line Summary for Examiners}\\[4pt]
\textit{``Fixed-slope 87L over-restrains when $\kibr > 0.5$ because $\rho_\mathrm{min} = 2(1-\kibr)/(1+\kibr) < \km$. SAMBP-87L tracks $\hat{\kibr}$ every 20~ms (Lyapunov-stable, $\dot{V} < 0$) and sets $\km = \kbase(1-\kibr)^2/(1+\kibr)^2 + \kfloor$. Three additional elements cover sequence-blind and travelling-wave scenarios. Combined: 12/21 (57\%) $\to$ 21/21 (100\%). Stage-2 gate ($\gamma \geq 0.85$) prevents false trips under CT saturation.''}
\end{tcolorbox}
\end{center}

\end{document}
